% !TEX root = ../jlt.tex
\section{Featherweight Jlt syntax}
\label{app:fjlt:syntax}

FJlt is a subset of \Cpp{23}, parametric in a \emph{signature} $\Sigma$ of carriers.
A carrier is in $\Sigma$ on three conditions: it is \cppinline{std::semiregular}; the operations a program uses on it are \cppinline{constexpr} and total; and it has an image in the target, a constant type with those operations as constants, and an iterator if it is inductive.
The compiler decides semiregularity and that an operation is \cppinline{constexpr}; totality is a reading, since neither a \cppinline{constexpr} declaration nor a successful evaluation shows an operation defined on every input, and it is what Lemma~\ref{lem:progress} needs and what keeps \cppinline{int}'s overflow and division out; the image is a definition.
A type enters $\Sigma$ by meeting the conditions, not by being named here, which is the structural admission the fragment is named for; the carriers this paper's listings use are \cppinline{bool}, \cppinline{size_t}, \cppinline{std::string}, with concatenation and the comparisons as its operations, allocation failure aside, and the finite enumerations a program or the library declares, such as \cppinline{Ternary} and \cppinline{ArrowLaw}, each a finite sum of units with equality.
Figure~\ref{fig:fjlt-syntax} gives the grammar in the style of Featherweight Java: $\overline{\alpha}$ abbreviates a sequence $\alpha_1, \ldots, \alpha_n$, a construct written over a sequence is repeated for each element, and the separators the host language wants (commas, \cppinline{&&}, semicolons) are implied.
The overbar is also how a parameter pack reads: a variadic declaration is the family of its arities, each a declaration of the grammar, and a fold expression at arity $n$ is the nested term; what the overbar does not give is a statement about all arities at once.
Every FJlt program is a Jlt program, a \Cpp{23} program that keeps the three rules, with the same meaning, provided no law and no \cppinline{static_assert} builds a box (Appendix~\ref{app:fjlt:evaluation}).
The converse, for a Jlt program written in this grammar, is the adequacy argument of Section~\ref{sec:translation:adequacy}.
The listings are Jlt programs, and FJlt programs up to the sugar of Appendix~\ref{app:fjlt:closure}, with \cppinline{static_assert(e)} read as the claim that the closed term $e$ evaluates to \cppinline{true}; where a listing steps outside the kernel the text says so: Listing~\ref{lst:rule3}'s Collatz chain is finite by fact, not by a condition on $\mathcal{R}$, and Listing~\ref{lst:sites} dereferences an optional in a constant initialiser (Appendix~\ref{app:fjlt:closure}).

Metavariables: $\alpha, \beta$ range over type variables, $b$ over the carriers of $\Sigma$, $O$ over the function-object types, the standard templates (\cppinline{std::plus}, \cppinline{std::logical_and}, \ldots) and the operations a program declares, $r$ over ranges, $\oplus$ over the operations of $\Sigma$, total by the second condition: on unsigned carriers addition, subtraction, multiplication and the bitwise operations, which wrap, the comparisons, and division, remainder and the shifts by a literal the carrier can take (a nonzero divisor, a shift below the width); on \cppinline{bool} the logical operations; by a variable, division, remainder and the shifts are partial and not admitted, and the standard function objects $O$ are those of the same operations, $C$ over concept names, the standard ones included (\cppinline{std::semiregular} is the shape Rule~1 requires of every carrier), $S$ over the two standard result-type concepts \cppinline{std::same_as} and \cppinline{std::convertible_to}, $f$ over function names, $x$ over term variables and $c$ over constants.

\paragraph{Reading the grammar.}
A program is a sequence of declarations and a closed term.
A shape is a concept; its requirements are nested shapes, laws (closed \cppinline{bool} terms) and operations with their result types, exact (\cppinline{std::same_as}) or up to conversion (\cppinline{std::convertible_to}).
A constraint may be negated, \cppinline{!C<T>}, as a \cppinline{requires} clause can write it; a negation carries no evidence, and at concrete types it holds where $C$ does not (\textsc{E-Not}).
A trait is a function template declared and deleted, indexed by a carrier and a value of the operation; an instance is an explicit specialisation of it at a carrier; an instance with a context is a generic function whose \cppinline{requires} clause is the context.
A generic function is a \cppinline{constexpr} function template with a \cppinline{requires} clause and a single \cppinline{return}; \cppinline{auto} in a listing names the type the body has, and a parameter written \cppinline{const τ} or \cppinline{const τ&} is the parameter $\tau\; x$ of the grammar, since a value is immutable and a reference to one denotes it.
The standard function objects and the operations a program declares (a struct with fields and a call operator, itself generic where the listing writes \cppinline{auto}; a former such as \cppinline{Compose} is one) are the only function objects; a term applies one by constructing it, with its fields, and reads a field with \cppinline{.}; the call operator is \cppinline{const}, so applying a function object observes it, as Rule~1 requires of a value.
Terms are variables, constants, function names, calls with explicit or deduced type arguments, applications, arithmetic and comparisons, conversions between carriers, conditionals, pairs, sums and boxed values.
A sum is a type of a family $\mathcal{S}$ that, like $\mathcal{R}$ below, is defined by conditions: two injections $\iota_1, \iota_2$ and a total case form, \texttt{case(}$s$\texttt{, }$f$\texttt{, }$g$\texttt{)} applying $f$ or $g$ to the injected value, each instance spelling them in \Cpp{}.
\cppinline{optional<τ>} is $1 + \tau$, with $\iota_1$ = \cppinline{nullopt}, $\iota_2\,e$ = \cppinline{optional\{e\}} and the case \cppinline{s.transform(g).value_or(f)}, $f$ the value for the empty one; \cppinline{variant<τ₁, τ₂>} is $\tau_1 + \tau_2$, with $\iota_i\,e$ = \cppinline{variant<τ₁, τ₂>(in_place_index<i-1>, e)}, both alternatives spelled, and the case dispatched on the index, \cppinline{s.index() == 0 ? f(get<0>(s)) : g(get<1>(s))}, each \cppinline{get} guarded by the test, as the library's \cppinline{Copairing} writes it; \cppinline{std::visit} is the special case in which one operation, its call operator entailed at both summands, serves as both branches.
\cppinline{*} on an optional and \cppinline{std::get} on a variant are partial and are not terms; under a guard, or in a constant initialiser, where constant evaluation refuses the empty case, they are the case form.
A box is built and read, and nothing else: there is no term for \cppinline{get}, \cppinline{use_count} or \cppinline{==} between boxes, so no program can ask which box holds a value, and that absence is what makes a box a value.
There are two recursion primitives, one per kind of finiteness.
\cppinline{fold_left} folds over a range of fixed length, a value $r$ of a family $\mathcal{R}$ that, like $\Sigma$, is defined by conditions and not by a list: $r$ has a length $|r| \in \mathbb{N}$, fixed before the fold starts, and, when $|r| > 0$, a head and a tail $\mathit{next}\,r$ with $|\mathit{next}\,r| = |r| - 1$; \textsc{S-Fold} consults $|r|$ before it asks for either, and on both sides both are extended to the empty range, the head by the type's default value and the tail by the range itself, which the iterator never reaches, since it runs $|r|$ times.
The default value $\mathit{default}_\tau$ of a type is given by its structure, not by \cppinline{T{}}: a carrier's own default, a pair of defaults, the empty optional, the first injection of a default, a fieldless function object, and for a box a box of its content's default, so that no term of the kernel holds a null box.
Each instance $R$ of $\mathcal{R}$ over an element type $\tau$ contributes to the target a state type $\llbracket R \rrbracket$, a constructor constant $c_R$ applied to the translations of its arguments, and its operations $|\cdot| : \llbracket R \rrbracket \to \llbracket \texttt{size\_t} \rrbracket$, $\mathit{head} : \llbracket R \rrbracket \to \llbracket \tau \rrbracket$ and $\mathit{next} : \llbracket R \rrbracket \to \llbracket R \rrbracket$, constants for a base range and terms over the parts for a range built on another; the state of \cppinline{n >> chain(prior)} is $\llbracket N \rrbracket$, the count, that of \cppinline{transform(r, g)} is $\llbracket R \rrbracket \times \llbracket g \rrbracket$, the arrow carried so that $\mathit{head}$ applies it and $\mathit{next}$ keeps it, and that of \cppinline{iota(i, j)} is $\mathbb{N} \times \mathbb{N}$.
The kernel's ranges are these.
\cppinline{n >> chain(prior)} is the chain of a count under Figure~\ref{fig:structures}'s $\mathit{prior}$.
Its head is \cppinline{prior(n)} and its tail \cppinline{prior(n) >> chain(prior)}.
It is empty where \cppinline{prior(n)} is nothing, and finite because the step descends (Section~\ref{sec:posture}).
\cppinline{transform(r, g)} is the image of a range under an arrow, Listing~\ref{lst:rule3}'s \cppinline{r >> g}.
Its head is \cppinline{g(head r)}, its tail \cppinline{transform(next r, g)}, and its length $|r|$.
\cppinline{iota(i, j)}, with $i \le j$, is the naturals from $i$ below $j$, with head $i$ and tail \cppinline{iota(i + 1, j)}.
The condition $i \le j$ is the standard view's: reversed endpoints are undefined behaviour, not an empty range.
A range's operations are the terms \texttt{size(}$r$\texttt{)}, \texttt{head(}$r$\texttt{)} and \texttt{next(}$r$\texttt{)}, written $|r|$, $\mathit{head}\,r$ and $\mathit{next}\,r$ in the rules, typed by \textsc{T-Range} at the range's own type $R\texttt{<}\tau\texttt{>}$ and total: on the empty range $\mathit{head}$ is the default value of $\tau$, which Rule~1's semiregularity supplies, and $\mathit{next}$ is the range itself (\textsc{S-Range}), values \textsc{S-Fold} never reaches; in \Cpp{} the length is \cppinline{ranges::size}, the head \cppinline{*ranges::begin} under a guard on the length, as \cppinline{*} on an optional is admitted (Appendix~\ref{app:fjlt:closure}), and the tail has no standard spelling at the same type, so it is the instance's own.
A range may also be a type variable $R$ under the shape $\mathit{IsRange}\texttt{<}R\texttt{, }T\texttt{>}$ $=$ \cppinline{requires (const R r) { { size(r) } -> same_as<size_t>; { head(r) } -> same_as<T>; { next(r) } -> same_as<R>; }}, which the instances of $\mathcal{R}$ over $\tau$, and only they, satisfy at $T = \tau$ by \textsc{E-Shape}, each operation a field and the family's conditions, that the length falls by one along $\mathit{next}$, the premise, so a type with the three operations and no such law does not satisfy it; Listing~\ref{lst:rule3}'s \cppinline{concat} writes the weaker standard constraint \cppinline{std::ranges::input_range<R>}, which the kernel reads as $\mathit{IsRange}$ with $T$ the deduced \cppinline{ranges::range_value_t<R>}, so that the listing is in the kernel at the kernel's own ranges, and a standard range outside $\mathcal{R}$, an unsized view, is not.
\cppinline{cata(alg, t)} folds over a term: an operation whose fields are terms is a node, a term is a finite composition of nodes, and the fold applies the algebra to each node and the folded values of its legs, one node at a time.
A pair's projections are written as the standard writes them, \cppinline{e.first} and \cppinline{e.second}, and read as $\pi_1$ and $\pi_2$~\cite{pierce1991basic}; the library's \cppinline{π_1} and \cppinline{π_2} are the same projections as function objects, applications under \textsc{T-Op}.
The type formers are productions, not members of $\Sigma$, because the translation acts on them structurally: pairs as products, optionals as sums, boxes as unary products; a former enters when it has such an image, as \cppinline{std::variant} does as a sum, above, and \cppinline{std::vector} would as lists with their fold, which is absent because no listing needs it.

\begin{figure}[htbp]
\footnotesize
\[
\begin{array}{lcl@{\hspace{0.6em}}l}
P &::=& \overline{D}\; e & \text{program} \\[3pt]
D &::=& \texttt{template <typename }\overline{\alpha}\texttt{> concept } C \texttt{ = } \overline{R}\texttt{;} & \text{shape} \\
  & | & \texttt{template <typename }\overline{\alpha}\texttt{> constexpr } \tau\; f\texttt{(}\overline{\tau}\texttt{) = delete;} & \text{trait} \\
  & | & \texttt{template <> constexpr } \tau\; f\texttt{<}\overline{\tau}\texttt{>(}\overline{\tau\; x}\texttt{) \{ return } e\texttt{; \}} & \text{instance} \\
  & | & \texttt{template <typename }\overline{\alpha}\texttt{> requires } \overline{Q}\texttt{ constexpr } \tau\; f\texttt{(}\overline{\tau\; x}\texttt{) \{ return } e\texttt{; \}} & \text{generic function} \\
  & | & \texttt{template <typename }\overline{\alpha}\texttt{> struct } O \texttt{ \{ } \overline{\tau\; y}\texttt{;} & \text{operation} \\
  &   & \qquad \texttt{template <typename }\overline{\beta}\texttt{> requires } \overline{Q} & \\
  &   & \qquad \texttt{constexpr } \tau \texttt{ operator()(}\overline{\tau\; x}\texttt{) const} & \\
  &   & \qquad \texttt{\{ return } e\texttt{; \} \};} & \\[3pt]
R &::=& Q & \text{nested shape} \\
  & | & e & \text{law} \\
  & | & \texttt{requires (}\overline{\texttt{const }\tau\; x}\texttt{) \{ } \overline{\texttt{\{ }e\texttt{ \} -> }S\texttt{<}\tau\texttt{>;}} \texttt{ \}} & \text{operations} \\[3pt]
Q &::=& C\texttt{<}\overline{\tau}\texttt{>} \;\mid\; \texttt{!}C\texttt{<}\overline{\tau}\texttt{>} & \text{constraint} \\[3pt]
\tau &::=& \alpha \;\mid\; b \;\mid\; O\texttt{<}\overline{\tau}\texttt{>} \;\mid\; \texttt{pair<}\tau\texttt{, }\tau\texttt{>} \;\mid\; \texttt{optional<}\tau\texttt{>} & \text{type} \\
  & \mid & \texttt{variant<}\tau\texttt{, }\tau\texttt{>} \;\mid\; \texttt{shared\_ptr<const }\tau\texttt{>} & \text{sum, boxed value} \\
  & \mid & \mu F \;\mid\; R\texttt{<}\tau\texttt{>} & \text{term type, the nodes that build it; a range of instance } R \text{ over } \tau \\[3pt]
e &::=& x \;\mid\; c \;\mid\; f \;\mid\; f\texttt{<}\overline{\tau}\texttt{>(}\overline{e}\texttt{)} \;\mid\; f\texttt{(}\overline{e}\texttt{)} \;\mid\; O\texttt{<}\overline{\tau}\texttt{>\{}\overline{e}\texttt{\}} \;\mid\; e\texttt{(}\overline{e}\texttt{)} \;\mid\; e\texttt{.}y & \text{term} \\
  & \mid & e \oplus e \;\mid\; \tau\texttt{(}e\texttt{)} \;\mid\; e \texttt{ ? } e \texttt{ : } e \;\mid\; \texttt{pair\{}e\texttt{, }e\texttt{\}} \;\mid\; e\texttt{.first} \;\mid\; e\texttt{.second} & \text{conversion between carriers} \\
  & \mid & \iota_1\,e \;\mid\; \iota_2\,e \;\mid\; \texttt{case(}e\texttt{, }e\texttt{, }e\texttt{)} & \text{sum} \\
  & \mid & \texttt{make\_shared<const }\tau\texttt{>(}e\texttt{)} \;\mid\; \texttt{*}e & \text{box, unbox} \\
  & \mid & \texttt{fold\_left(}r\texttt{, }e\texttt{, }e\texttt{)} \;\mid\; \texttt{size(}r\texttt{)} \;\mid\; \texttt{head(}r\texttt{)} \;\mid\; \texttt{next(}r\texttt{)} & \text{fold over a range; its length, head and tail} \\
  & \mid & \texttt{cata(}e\texttt{, }e\texttt{)} & \text{fold over a term} \\[3pt]
r &::=& \text{a range of } \mathcal{R}\text{, with a length } |r| \text{ and, when } |r| > 0\text{, a head and a tail} & \text{range of fixed length} \\
\end{array}
\]
\normalsize
\caption{The syntax of FJlt, a subset of \Cpp{23}.
Declarations, requirements, constraints, types, terms and ranges, namespace qualifiers elided; the paragraph \emph{Reading the grammar} in Appendix~\ref{app:fjlt:syntax} reads it line by line.}
\label{fig:fjlt-syntax}
\end{figure}

